\section{Reciprocal inversion without a common drift direction}
\label{sec:mean-exit}
The finite stopping hypothesis of Theorem~\ref{thm:stopped} is sufficient
but is not necessary for identification. A centered random displacement
can remain inside one component for arbitrarily many steps. Nevertheless,
its mean exit time is uniformly bounded by the component diameter and
its second moment. We use that bound twice: to recover the occupation
without specifying its zero set, and to compare two unknown occupations
with a constant independent of the iteration depth.

\subsection{The observed forcing and the computational walk}
Let $\rho$ be a prescribed probability measure on $\R^d$. The forward
protocol draws $q\sim f$ and $a\sim\rho$ independently and attempts the
straight segment from $q$ to $q+a$. The reverse protocol draws the same
joint law for $(q,a)$ but attempts the segment from $q+a$ to $q$.
Each protocol records only its hit bit. A solid start and a free miss
are both zero; all attempts are counted. No displacement-resolved
subpopulation mean is required. The direction and duration randomizer
is part of the input controller, not an additional geometric output.

With $\chi=\one_{\mathcal O}$ and
$T_\rho w(x)=\int w(x+a)\rho(\dd a)$, the pointwise reversal balance
and Fubini give
\begin{equation}\label{eq:reciprocal-forcing}
 P_\rho^+(f)-P_\rho^-(f)=\int f(T_\rho-I)\chi.
\end{equation}
For the translated even kernel commands of Section~\ref{sec:acquisition},
translation invariance commutes with convolution. Thus the observed
balance function is $g_\sigma=(T_\rho-I)u_\sigma$, where
$u_\sigma=\kappa_\sigma*\chi$. The corresponding random walk is a
\emph{computational} process. Its successive locations are not a
multi-collision trajectory and its exit from an unknown component is
not an observed stopping signal.

In this section impose, for known $b,v_*>0$,
\begin{equation}\label{eq:centered-law}
 |a|\le b\quad\rho\text{-a.s.},\qquad
 \int a\,\rho(\dd a)=0,\qquad
 s_2:=\int|a|^2\rho(\dd a)\ge v_*.
\end{equation}
There need not be a direction with positive drift. Neither symmetry,
a density for $\rho$, nor a positive definite covariance matrix is
required. The full angular circle, or two equally weighted opposite
steps, satisfies these conditions when the nonzero step length is fixed.
The reverse translation in \eqref{eq:reciprocal-forcing} remains essential;
independently choosing a reverse position with the same marginal
angular law is a different experiment.

\begin{lemma}[A uniform mean exit bound]\label{lem:mean-exit}
Let $0\le u\le1$ be measurable. Suppose its positive set is contained
in a union of sets of diameter at most $D$, separated from one another
by a distance strictly greater than $b$. Let $X_n=x+a_1+\cdots+a_n$,
where the increments are independent with law \eqref{eq:centered-law},
and put $\tau_u=\inf\{n\ge0:u(X_n)=0\}$. Then
\begin{equation}\label{eq:exit-bound}
 \sup_x\E_x\tau_u\le H:=\frac{(D+b)^2}{v_*},\qquad
 \sup_x\Prob_x(\tau_u>n)\le2^{-\lfloor n/L\rfloor},
 \quad L=\max(1,\lceil2H\rceil).
\end{equation}
No regularity of the positive set is used.
\end{lemma}
\begin{proof}
If $u(x)=0$, the exit time is zero. Otherwise, until that time all
successive locations with positive $u$ lie in the same containing set:
a step of length at most $b$ cannot jump to a different one. For
$S=\tau_u\wedge n$ we consequently have $|X_S-x|\le D+b$.
Independence and zero mean make
$|X_n-x|^2-ns_2$ a martingale. Applying the bounded stopping identity
at $S$ gives
\[
 s_2\E_x S=\E_x|X_S-x|^2\le(D+b)^2.
\]
Monotone convergence proves the expectation estimate, including
almost-sure finiteness of $\tau_u$. Markov's inequality gives
$\Prob_x(\tau_u>L)\le H/L\le1/2$. At time $L$, on survival,
the same expectation estimate applies from the current location.
Iteration of the Markov property gives the stated geometric tail.
This argument does not presume optional stopping at an unbounded time;
only bounded stopping is used before taking the limit.
\end{proof}

\begin{theorem}[Mean-exit inverse and physical comparison]
\label{thm:mean-exit-inverse}
Under Lemma~\ref{lem:mean-exit}, write $g=(T_\rho-I)u$ and define
\begin{equation}\label{eq:mean-exit-recursion}
 V_0=0,\qquad V_{n+1}=\max\{0,T_\rho V_n-g\}.
\end{equation}
Then $0\le V_n\le V_{n+1}\le u$ and
\begin{equation}\label{eq:mean-exit-rate}
 \|u-V_n\|_\infty\le2^{-\lfloor n/L\rfloor}.
\end{equation}
If $\widehat g$ has uniform error at most $\xi$, and each transition
update has additional uniform error at most $\zeta$, clipping the
perturbed iterates to $[0,1]$ gives
\begin{equation}\label{eq:mean-exit-noise}
 \|u-\widehat V_n\|_\infty
       \le2^{-\lfloor n/L\rfloor}+n(\xi+\zeta).
\end{equation}
In particular, increasing the iteration depth without regard to the
error is not the asserted regularization.

There is a stronger comparison between physical candidates. Suppose
$u,\widetilde u\in[0,1]$ both satisfy the same containing-set bounds,
with no prescribed agreement of their zero sets. If
$g=(T_\rho-I)u$ and $\widetilde g=(T_\rho-I)\widetilde u$, then
\begin{equation}\label{eq:unknown-support-stability}
 \|u-\widetilde u\|_\infty
                  \le H\|g-\widetilde g\|_\infty.
\end{equation}
Thus $g$ determines $u$ in this class. This estimate is independent of
both a convolution scale and the iteration depth. It compares admissible
occupations, not arbitrary noisy outputs of the recursion.
\end{theorem}
\begin{proof}
Positivity of $T_\rho$ and $T_\rho u-g=u$ give the monotone bounds.
Every iterate vanishes where $u=0$. For $e_n=u-V_n$, at a point with
$u>0$ the recursion gives $e_{n+1}=\min(u,T_\rho e_n)\le T_\rho e_n$.
Iterating until the first zero yields
\[
 0\le e_n(x)\le\E_x\{u(X_n)\one_{\{\tau_u>n\}}\}.
\]
Lemma~\ref{lem:mean-exit} proves \eqref{eq:mean-exit-rate}.
The maximum, clipping and a Markov transition are nonexpansive in
supremum norm. Induction proves \eqref{eq:mean-exit-noise}.

For the comparison put $w=u-\widetilde u$ and
$\epsilon=\|g-\widetilde g\|_\infty$. Bounded stopping in the
martingale for $w$ gives, with $S=\tau_u\wedge n$,
\[
 w(x)=\E_x w(X_S)
       -\E_x\sum_{j<S}(g-\widetilde g)(X_j).
\]
At $\tau_u$ the terminal difference is $-\widetilde u\le0$; on
survival its positive part is at most one. Hence
$w(x)\le\Prob_x(\tau_u>n)+\epsilon\E_x(\tau_u\wedge n)$.
Letting $n$ increase gives $w(x)\le H\epsilon$. Interchanging the
two functions and stopping at $\tau_{\widetilde u}$ proves the reverse
bound. This uses neither an observed free reference point nor knowledge
of either stopping set in the reconstruction.
\end{proof}

\begin{corollary}[Every nontrivial bounded displacement law]
\label{cor:noncentered-exit}
The same identification and fixed-aperture comparison hold for any
prescribed law supported in $|a|\le b$ other than $\delta_0$, provided
the containing sets have separation greater than $b$. More explicitly,
suppose for some unit vector $v$ and numbers $\gamma,p>0$ that
$\rho\{a:a\cdot v\ge\gamma\}\ge p$. Choose an integer $m$ with
$m\gamma>D$. Then
\[
 \Prob_x(\tau_u>n)\le(1-p^m)^{\lfloor n/m\rfloor},\qquad
 \E_x\tau_u\le m/p^m.
\]
In Theorem~\ref{thm:mean-exit-inverse} and
Proposition~\ref{prop:fixed-exit-aperture}, replace the survival bound
and $H$ by these quantities. The centered-law bound
\eqref{eq:exit-bound} is usually substantially better; the constants
here can be large when the escape cone has small probability.
\end{corollary}
\begin{proof}
Conditional on any surviving starting position, the next $m$
independent increments all belong to the displayed cone with probability
at least $p^m$. If the occupation stayed positive throughout those
steps, they would lie in one containing set and their net projection
would exceed its diameter. Hence this event forces a visit to zero.
The Markov property at block endpoints gives the tail. Summing the tail
in blocks of length $m$ gives the expectation bound. The preceding
Bellman and stopped-comparison proofs require only these two bounds.
Finally, any law different from $\delta_0$ has a nonzero point in its
support; a sufficiently small ball about that point lies in a cone
$a\cdot v\ge\gamma>0$ and has positive probability. Thus the stated
condition does not require an almost-sure common drift direction.
\end{proof}

For mollified occupations under the physical component bounds, take
\begin{equation}\label{eq:uniform-exit-prior}
 D=D_0+2\sigma_0,\qquad b<d_0-2\sigma_0,
 \qquad 0<\sigma\le\sigma_0.
\end{equation}
Their positive sets lie in the $\sigma$-enlarged components, which
satisfy these diameter and separation bounds. Thus the same $H,L$
work at every scale. Equality of the exact two functionals on smooth
spatial commands determines every $g_\sigma$, hence every $u_\sigma$.
Approximate-identity convergence determines $\chi$ almost everywhere.
For the closed convex physical components this determines the sets
as closures of their interiors. Analyticity and periodicity are not
needed in this local inference.

\subsection{Locality, stability of the prescribed law, and sharpness}
An infinite iteration must not silently turn a local experiment into
an unbounded aperture experiment. The following localization resolves
that issue.

\begin{proposition}[A fixed aperture for the mean-exit inverse]
\label{prop:fixed-exit-aperture}
Let $U$ be a bounded target set and let a bounded observation set $W$
contain $U+\overline B_{D+b+\varepsilon_0}$ for some
$\varepsilon_0>0$. Define $T^Wv=T_\rho(\one_Wv)$ and perform
\eqref{eq:mean-exit-recursion} on $W$, setting every iterate to zero
outside $W$. Then \eqref{eq:mean-exit-rate} and
\eqref{eq:mean-exit-noise} hold on $U$, using only forcing errors on
$W$. The comparison \eqref{eq:unknown-support-stability} on $U$ likewise
requires forcing data only on this fixed enlarged set.
\end{proposition}
\begin{proof}
Globally the zero extension gives $T^Wu\le T_\rho u$, so induction
still gives $0\le V_n\le u$ and zero values on $\{u=0\}$.
For a start $x\in U$ with $u(x)>0$, the entire containing set of $x$
and every possible exit location lie in $W$. Along the walk up to
its first zero, $T^W$ therefore introduces no boundary discrepancy.
The killed error and stopped comparison proofs apply verbatim along
this walk. Errors elsewhere in $W$ can propagate through at most $n$
updates, giving the same bound as before. The argument on $U$ does
not require the unknown occupation to vanish near $\partial W$.
\end{proof}

Define total variation by
$\|\rho-\rho'\|_{\rm TV}=\sup_{0\le h\le1}|\int h\,\dd(\rho-\rho')|$.
For a physical $u\in[0,1]$,
\begin{equation}\label{eq:operator-calibration}
 \|(T_\rho-T_{\rho'})u\|_\infty
                    \le\|\rho-\rho'\|_{\rm TV}.
\end{equation}
Consequently an independently certified error in a reciprocal randomized
law can be charged directly to the forcing error. A baseline law
satisfying \eqref{eq:centered-law} suffices for this perturbation
comparison; the small error need not itself have mean zero. The actual
forward and reverse preparations must still be reciprocal. Uncoupled
position, timing and angle errors are instead charged by
Theorem~\ref{thm:calibration}, with its $\sigma^{-1}$ cost.
We do not infer either calibration guarantee from the observed bits.

The expectation factor has a familiar potential-theoretic meaning. If
$Q$ is a finite killed transition matrix with finite mean exit times,
then $G=(I-Q)^{-1}=\sum_{n\ge0}Q^n$ has nonnegative entries and
\begin{equation}\label{eq:green-norm}
             \|G\|_{\infty\to\infty}=\max_x\E_x\tau.
\end{equation}
Indeed its row sums are the sums of survival probabilities. The constant
is attained by forcing identically one on the killed state set. This is
a sharp Dirichlet resolvent norm, not a minimax lower bound for the
physical collision experiment with unknown obstacles. Green functions
and stopped martingales are classical tools; see \cite{LawlerLimic}.
Our use of them concerns a measured forcing and an unknown zero set,
with the containing sets needed only for bounds and not given as data.

For the symmetric nearest-neighbor walk on the line and $u=1$ at two
adjacent sites and zero elsewhere, \eqref{eq:mean-exit-recursion} gives
$V_n=1-2^{-n}$ at those sites. The mean exit time is two, although no
finite almost-sure stopping bound exists. This example belongs to the
new theorem and verifies why it is a genuine extension of
Theorem~\ref{thm:stopped}, not a claim that its earlier finite stopping
obstruction has disappeared. If the second moment vanishes, constants
on a stationary orbit cannot be recovered; the positive moment in
\eqref{eq:centered-law} has a different role from a common positive drift.
